\section{Runtime Components}

Compilation and instrumentation equip the target binary with the tools to report execution paths, but once visibility into the target is established, the challenge of CGF shifts to execution throughput. Processing millions of mutated inputs via standard system execution flows can introduce heavy performance bottlenecks.

The fuzzer requires an execution environment capable of managing the fuzzing loop, delivering mutated inputs, and evaluating the target's behaviour at scale. To achieve efficiency, modern fuzzers use high-speed inter-process communication (IPC) and specialized state management.

\subsection{The Forkserver Architecture}

The normal system execution flow for executing a program is to call \texttt{fork} and \texttt{execve}, but the latter is extremely slow, as it requires the kernel to load the binary, set up the memory layout, and invoke the dynamic linker. 

The \textit{forkserver} avoids this overhead. The fuzzer executes the target program only once. The target program initializes itself and stops before the main fuzzing loop, either automatically at \texttt{main} or via \texttt{\_\_AFL\_INIT()}. At this point, it acts as a "forkserver": for every fuzzing iteration, the fuzzer sends a signal to the forkserver, which simply calls \texttt{fork}, letting the newly created child process inherit the already-initialized memory state. The child then executes the target function with the new input and exits. The forkserver parent waits for the child and reports its exit status back to the fuzzer.

\paragraph{Communication Channels} Fuzzer and forkserver communicate via two standard pipes: 
\begin{itemize}
    \item A \textit{Control} pipe \texttt{FORKSRV\_FD}, used by the fuzzer to send commands to the forkserver.

    \item A \textit{Status} pipe \texttt{FORKSRV\_FD + 1}, used by the forkserver to send status updates back to the fuzzer.
\end{itemize}

\paragraph{Handshake protocol} AFL++ introduces a protocol to allow fuzzer and target to negotiate capabilities dynamically at startup.

It includes an initial handshake to check the supported forkserver version; then, the fuzzer and target negotiate advanced features through a 32-bit bitmask.

\paragraph{Execution Loop} Once the handshake is complete, the standard execution loop begins for each input generated by the fuzzer: 
\begin{enumerate}
    \item The fuzzer writes a 4-byte command to the control pipe, typically instructing it to start.

    \item The forkserver reads the command. If the pipe is closed, the fuzzer exited and the forkserver terminates itself. Otherwise, it calls \texttt{fork}, lets the child run the target logic with the new input, and writes the PID of the child process to the status pipe.

    \item The forkserver then waits for the child to finish and writes the exit code to the status pipe.

    \item The fuzzer reads the status, updates the coverage bitmap, calculates the feedback and repeats the cycle.
\end{enumerate}

\paragraph{Alternative Backends} The forkserver architecture allows various execution backends, such as the \textit{fauxserver} --- which uses \texttt{execv} for targets incompatible with the fork architecture --- or other modes such as the ones described in Section \ref{subsec:bbb-instr}.

\subsubsection{Persistent Mode}

Even with a forkserver, calling \texttt{fork} thousands of times per second incurs significant overhead. Persistent mode solves this by doing a single \texttt{fork} and inserting a loop around the target function. The child then repeatedly reads the fuzzer input, executes the target function and resets any necessary state, without ever exiting. 

When the loop counter reaches its limit, the child exits and a new one is created from the forkserver. This limits the impact of memory leaks while still achieving superior performance.

AFL++ includes features for tracking persistent mode executions. Because the target state might not be perfectly reset between iterations in persistent mode, a crash might depend on the specific sequence of inputs executed before the one that causes the crash. With \texttt{AFL\_PERSISTENT\_RECORD} enabled, the forkserver tracks the last $n$ inputs processed by the child, allowing the fuzzer to replay the exact state transitions that led to the crash.

\subsection{Shared Memory}
\label{sub:shm}

The primary mechanism used to orchestrate high-throughput communication is shared memory. By mapping common memory regions into the address spaces of both fuzzer and target, AFL++ avoids the latency overhead associated with traditional IPC --- such as using pipes or file system I/O.

The shared memory architecture in AFL++ relies on zero-copy data exchange. When the fuzzer is initiated, it allocates memory regions for feedback data and operational purposes. Upon execution, the target binary retrieves, via environment variables, the identifiers for these pre-allocated regions, attaching them to its own virtual address space.

\subsubsection{Lifecycle and Allocation Mechanisms}

The management of the shared memory is mainly governed by a single module (\texttt{afl-sharedmem.c}) which abstracts the OS paradigms into an API, defined in the corresponding header file. There is a data structure maintaining everything related to the state, file descriptors, object identifiers and virtual memory pointers for the shared regions.

The initialisation routine handles allocating and preparing the necessary memory segments, as well as constructing the corresponding environment strings. Along with the primary coverage map, it handles the creation of auxiliary maps, when enabled (e.g., CmpLog, ASan). A deinit function provides the cleanup counterpart, unmapping virtual memory and unsetting all environment variables.

\paragraph{Large Page Allocation} To minimize Translation Lookaside Buffer (TLB) misses during execution loops, AFL++ attempts to utilize \textit{large pages}. This decreases memory management overhead and increases performance.

\subsubsection{Coverage Map}

The primary coverage map acts as the central data structure for evaluating execution feedback. Instantiated within the shared memory region during the initialisation phase, it is conventionally structured as a fixed-size (although configurable) contiguous byte array.

After the forkserver signals the termination of a child process, the fuzzer directly analyses the modified map without further IPC. Each byte in the array corresponds to the execution frequency of a specific control-flow edge. When a transition from basic block $A$ to $B$ occurs, the instrumentation logs this edge by incrementing a byte in the map.

The coverage map's lifecycle is split between two processes: the fuzzer --- which creates and analyses it --- and the target --- which attaches to and mutates it.

\paragraph{Post-Execution Analysis} After the target child process terminates, control returns to the parent process, which must efficiently analyse the resulting map to determine if the execution yielded interesting behaviour. 

The fuzzer first checks for novelty in the generated map by comparing it against a persistent \texttt{virgin\_bits} map, which tracks all previously-discovered edges and is updated when novel transitions are discovered.  

To combat path explosion caused by trivial variations in the number of loop iterations (e.g., observing an edge 7 times instead of 6), the fuzzer groups execution counts into one of 8 pre-defined buckets. Changes within the same bucket are ignored, mitigating state explosion.

The input is considered "interesting" if it returns new coverage or a new count for one of the buckets. The input is subsequently saved to the queue directory and becomes a seed for future mutation cycles.

\subsection{Queue and Scheduling}
\label{subsec:qandsched}

In AFL++, the queue is not a simple data structure. It serves as the evolving repository for the fuzzer's seed corpus. Every input that triggers new coverage or unique behaviour is preserved as a distinct queue entry, serving as material for future mutations. It is the queue's responsibility to orchestrate the fuzzing loop by deciding which entry to schedule next and how much \textit{energy} (i.e., fuzzing time) to allocate to it.

Before being actually inducted into the queue, the input undergoes a calibration phase. The test is executed multiple times to ensure the observed behaviour is deterministic rather than an anomaly.

Furthermore, there is an array that keeps the best queue entry --- in terms of speed and size --- for each byte of the coverage map. Only the best entry for each edge is kept.

\paragraph{Scoring} When selecting a test case from the queue to undergo mutation, AFL++ must determine the intensity and duration of the fuzzing cycle. This is codified as the \texttt{perf\_score}. The calculation of the score is a composition of several heuristics, all modulated by a \textit{power schedule}.

The base metrics taken into account are: 
\begin{itemize}
    \item \textit{Execution speed:} Faster inputs are prioritized, as they're easier to fuzz.

    \item \textit{Bitmap density:} Entries that traverse more unique edges --- i.e., have a fuller bitmap --- are assumed to be more valuable.

    \item \textit{Temporal handicap:} When a new queue entry is discovered late in the fuzzing campaign, it has inherently missed the previous mutation cycles that other entries have enjoyed; to mitigate this, newly discovered entries have a higher multiplier in scoring.

    \item \textit{Generational Depth:} The distance, in number of generations, from the starting corpus is treated as a proxy for semantic complexity; inputs with higher depth are given higher multipliers.
\end{itemize}

There are several power schedules, used as "macro-strategies" that fundamentally alter the calculated score, in one way or another.

\paragraph{Selection} AFL++ does not simply iterate over the queue sequentially. It uses weighted random selection via an alias table.

Each entry's weight is derived from its \texttt{perf\_score} and the next queue entry is selected by using the alias table. This ensures that \texttt{favored} and highly scored entries are selected far more frequently than redundant or slow entries.

\paragraph{Culling} To prevent the fuzzer from wasting time on redundant paths, a greedy algorithm is called periodically to calculate a minimal set of test cases that hit every known edge. These are marked as \texttt{favored}. During the main fuzzing loop, there is an extreme bias towards the \texttt{favored} entries.

Efficient entries often cover large portions of code; selecting a single entry from all the top-rated ones might clear thousands of bits from the queue. This ensures that the final set of favoured entries remains much smaller than the total queue.

\subsection{Compiler Runtime}

The instrumentation logic cannot work without the AFL++ compiler runtime (\texttt{afl-compiler-rt.o.c}). The runtime is responsible for bridging the gap between the instrumented edges embedded in the target binary and the fuzzer process. It is linked as a static object, and as such constructors inside the runtime are executed before the target main function is invoked.

The instrumentation process produces code that references symbols provided by the runtime, including generic variables such as the ones used for the shared memory.

The runtime operates primarily as an initialisation bootstrap and execution harness. It intercepts the normal execution flow of the target program to establish the shared memory region, initialise the forkserver, manage persistent mode, and support advanced features.

\paragraph{Shared Memory} AFL++ introduced dynamic map sizing to support larger targets without excessive map collisions. 

The instrumentation passes export a symbol representing the total number of instrumented edges. The runtime evaluates the value and, if larger than the default size, the runtime sets the size of the map to accommodate the exact number of edges.

Depending on the platform and configuration, the runtime attaches the shared memory using either POSIX shared memory or System V shared memory. The base address of the mapped memory is assigned to the global variable \texttt{\_\_afl\_area\_ptr}.

\paragraph{Forkserver} Inside the runtime resides the forkserver initialization hook, used to pause the process and establish the communication channel with the fuzzer process. 

Following the communication handshake, instead of terminating after a single input, the forkserver enters an infinite loop of:
\begin{itemize}
    \item Waiting for a command from the fuzzer.

    \item Cloning itself, executing program logic in the child.

    \item Waiting for the end of execution and reporting the status to the fuzzer process.
\end{itemize}

\paragraph{Deferred Initialisation} By default, AFL++ stops the forkserver just before the main function of the target process. While safe, this means that every cloned child must re-execute the entirety of the initialisation code. If a target performs expensive one-time setups --- e.g., parsing large configuration files --- fuzzing performance will be severely degraded.

Deferred initialisation solves this by allowing the forkserver to start at a point chosen by the user, where the expensive setup is completed.

In this case, the compiler runtime skips calling the forkserver constructor and instead calls the manual forkserver initialisation at the specified time. 

\subsection{Mutation Engine}

AFL++'s mutation engine is the component primarily responsible for input generation. Unlike blind random fuzzing, AFL++ employs a more sophisticated, coverage-guided evolutionary algorithm. It utilizes a structured and heuristic-driven pipeline to alter seed inputs. By continually adapting its strategies based on execution tracing, the engine ensures the computational effort is focused on generating the best possible inputs.

\paragraph{Phases} The standard execution of the default mutation engine for any given input follows a sequential progression:
\begin{enumerate}
    \item Calibration and Trimming: Before mutation, the input may be calibrated to confirm the execution trace and trimmed to reduce size without losing coverage.

    \item Deterministic stages: if enabled and the input has not already been subjected to them, AFL++ applies the following stages: bitflips, general arithmetic, insertion of common interesting values and dictionary insertions.

    \item Non-Deterministic Stage (Havoc): multiple mutations stacked on top of each other in a highly randomized fashion. The number of cycles depends on the seed's performance score and queue depth; favourable inputs receive significantly more havoc time than uninteresting ones.

    \item Splicing: if the havoc stage fails, the engine will attempt splicing an input with another one from the queue.
\end{enumerate}

The early stages guarantee that every single byte will be modified sequentially. This is guaranteed to find simple syntactic branches but can be computationally expensive. For these reasons, it is never repeated for the same queue entry once completed.